\begin{abstract}
A bank's support desk does not need a label. It needs to know which ticket to
open next. Published triage systems predict a per-ticket severity and report
classification metrics, leaving the ordering question unanswered. We present a
ticket ordering system for a trilingual Sri Lankan retail bank, where customers
write in Sinhala, Tamil, and English and render the two local languages in Latin
characters as readily as in their native scripts. The Ticket Urgency Score
combines a priority posterior, a sentiment posterior, and a per-intent
criticality prior estimated from training data, under multiplicative aging, with
weights fitted on development data and applied unchanged to test. Evaluated in a
discrete-event queue over 200 seeds, it attains 97--99\% of the reduction in
urgent-ticket waiting time that a perfect classifier would deliver.

The result that matters for deployment is a fairness one. Under a naive
argmax-tier policy, a customer who writes an urgent ticket in romanized Sinhala
waits \textbf{4.07 minutes longer} for a first response than one writing the
same ticket in English (95\% CI $[2.96, 5.21]$), while the native-script tracks
are indistinguishable from English. The comparison is controlled, the romanized
track being a deterministic transliteration of the native one. The gap is absent
under first-come-first-served and under a gold-label control, so the model causes
it rather than the arrival process, and our score cuts it to $+0.37$ minutes. We also show that classifier accuracy does not
transfer to the queue: a 2-point macro-F\textsubscript{1} advantage buys nothing
through the median and costs six hours in the tail, because the more accurate
model makes fewer errors but $8.5\times$ more confident ones.
\end{abstract}
